% Generated draft: requires mathematical and source review.
\documentclass[11pt]{article}
\usepackage[margin=25mm]{geometry}
\usepackage{fontspec}
\setmainfont{Latin Modern Roman}
\usepackage{amsmath,amssymb,mathtools,unicode-math}
\setmathfont{Latin Modern Math}
\usepackage{enumitem,needspace,xurl,hyperref}
\setlength{\parindent}{0pt}
\setlength{\parskip}{5pt}
\setlength{\emergencystretch}{3em}
\newcommand{\N}{\mathbb{N}}
\newcommand{\Z}{\mathbb{Z}}
\newcommand{\Q}{\mathbb{Q}}
\newcommand{\R}{\mathbb{R}}
\newcommand{\C}{\mathbb{C}}
\newcommand{\F}{\mathbb{F}}
\newcommand{\A}{\mathbb{A}}
\newcommand{\PP}{\mathbb{P}}
\newcommand{\E}{\mathbb{E}}
\newcommand{\eps}{\varepsilon}
\newcommand{\dd}{\,\mathrm{d}}
\begin{document}
\section*{UnsolvedMath ID: 701198}
\textbf{Source problem code:} catalogue entry
\subsection{Definitions and mathematical statement}
Let \(p\) be a prime. A group \(G\) is a locally finite \(p\)-group if every finitely generated subgroup of \(G\) is a finite group of order a power of \(p\). Write \(G'=[G,G]\) for its commutator subgroup; \(G\) is perfect if \(G=G'\). A subset \(S\subseteq G\) has bounded exponent if there is an integer \(n\geq 1\) such that \(s^{p^n}=1\) for every \(s\in S\). The requirement that \(S\) generate \(G\) means \(\langle S\rangle=G\); no finiteness assumption on \(S\) is imposed.

For a group \(H\), define its transfinite upper central series by
\[
Z_0(H)=1,\qquad Z_{\alpha+1}(H)/Z_\alpha(H)=Z\bigl(H/Z_\alpha(H)\bigr),\qquad
Z_\lambda(H)=\bigcup_{\alpha<\lambda}Z_\alpha(H)
\]
for every limit ordinal \(\lambda\). The group \(H\) is hypercentral if \(Z_\alpha(H)=H\) for some ordinal \(\alpha\). Define the derived series by \(H^{(0)}=H\) and \(H^{(k+1)}=[H^{(k)},H^{(k)}]\). The group \(H\) is solvable if \(H^{(d)}=1\) for some integer \(d\geq 0\).

\textbf{Question (A. O. Asar, Kourovka Notebook 16.6).} Can there exist a perfect locally finite \(p\)-group \(G\), a subset \(S\subseteq G\), and an integer \(n\geq 1\) such that
\[
\langle S\rangle=G,\qquad (\forall s\in S)\ s^{p^n}=1,
\]
under each of the following two separate additional hypotheses?
\begin{enumerate}
\item[(a)] Every proper subgroup \(H<G\) is hypercentral.
\item[(b)] Every proper subgroup \(H<G\) is solvable.
\end{enumerate}
The ordinal in (a) and the derived length in (b) may depend on \(H\); no uniform bound is specified. The exponent bound applies to the generating subset, not necessarily to every element of \(G\).

\textbf{Unresolved qualifications.} The supplied statement does not explicitly require \(G\neq 1\), and no such condition has been inserted here. Whether a nontrivial group is intended requires clarification. The source also does not specify whether the existence question is to be answered for each prime \(p\), or whether existence for some prime is the intended target.
\subsection{Short English statement}
The problem asks whether a perfect locally finite \(p\)-group can have generators of uniformly bounded order when every proper subgroup is hypercentral, and separately when every proper subgroup is solvable. The intended nontriviality requirement and quantification over primes are not explicit in the supplied source.
\subsection{Sources}
\begin{itemize}
\item Supplied local UnsolvedMath registry record: permanent ID \(701198\), exact source code \texttt{catalogue entry}, title \emph{Kourovka Notebook Problem 16.6}; upstream revision \texttt{d371ee917ac9647fbb06cd6fc529deec6f3dda82}. The registry labels the problem as open; that label is supplied metadata, not an independently verified mathematical assessment.
\item Attributed by the supplied record to A. O. Asar, proposed in 2006, in \emph{The Kourovka Notebook (Unsolved Problems in Group Theory)}, 16th issue (2006), problem 16.6. The record also identifies the 21st issue edition (2026), \texttt{arXiv:1401.0300v48}. Its collection description names E. I. Khukhro and V. D. Mazurov as editors. These bibliographic details are supplied attribution and have not been independently checked.
\end{itemize}
\end{document}
